% Guia do professor — Capítulo 16: Ciclones Tropicais
\section{Capítulo 16 --- Ciclones Tropicais}

\subsection*{Visão geral e ritmo}

O anti-Cap.~13: núcleo quente, sem frentes, movido a calor latente do
oceano. A joia pedagógica é o furacão como máquina de Carnot — um
motor térmico desenhado no céu que a turma de física reconhece na
hora. Fechar com Catarina (2004) traz o tema para casa.
\textbf{Ritmo: 2 aulas de 2\,h.}

\begin{itemize}
  \item \textbf{Aula 1} — anatomia (olho, parede, bandas), tabela de
        contraste com o Cap.~13, condições de gênese, Carnot/MPI e
        WISHE. Casos~1 e 3 (o \texttt{tcpyPI} impressiona).
  \item \textbf{Aula 2} — Holland e vento gradiente no limite,
        trajetórias e deriva beta, perigos com ênfase em surge
        (Caso~4), Catarina e o clima futuro (N3).
\end{itemize}

\subsection*{Objetivos de aprendizagem}

\begin{enumerate}
  \item contrastar ciclones tropicais e extratropicais (energia,
        núcleo, estrutura);
  \item enunciar as condições de gênese e usá-las para explicar a
        climatologia das bacias;
  \item calcular a intensidade potencial (Carnot simplificado e
        \texttt{tcpyPI}) e explicar o WISHE;
  \item obter o perfil de vento do modelo de Holland;
  \item estimar maré de tempestade (empilhamento $+$ barômetro
        invertido) e ranquear os perigos.
\end{enumerate}

\subsection*{Pré-requisitos a reativar}

Clausius--Clapeyron e $\theta_e$ (Cap.~4); fluxos bulk (Cap.~3);
vento gradiente e Coriolis (Cap.~10); hipsométrica para o núcleo
quente (Cap.~1); alísios e $\beta$ (Cap.~11); CAPE (Cap.~14).

\subsection*{Armadilhas conceituais frequentes}

\begin{itemize}
  \item \textbf{``Furacão é tempestade grande''}: é um sistema
        AXISSIMÉTRICO de núcleo quente — outra classe de objeto; as
        tempestades do Cap.~14 são os tijolos da parede.
  \item \textbf{Cisalhamento}: no Cap.~14 organiza; aqui MATA
        (ventila o núcleo quente). Vale uma linha no quadro.
  \item \textbf{``26,5\,°C é lei''}: é limiar empírico do $\Delta k$;
        com topo muito frio, menos serve (subtropicais, medicanes,
        Catarina).
  \item \textbf{Olho calmo $=$ fim}: a parede oposta chega em
        minutos — erro clássico (e fatal) de quem sai do abrigo.
  \item \textbf{Vento como maior perigo}: a água (surge $+$ chuva)
        mata mais; dano eólico $\propto v^3$, mas enchente não segue
        categoria.
\end{itemize}

\subsection*{Demonstração de baixo custo}

Ciclo de Carnot desenhado sobre o corte do furacão (transparência ou
quadro em duas cores) — a analogia motor térmico rende. Imagens de
satélite do Catarina (2004) e discussão: por que surpreendeu?
(Atlântico Sul ``não tinha'' furacões.)

\subsection*{Problemas sugeridos do livro (exercícios do Cap.~16)}

Aquecimento: MPI para TSMs dadas; queda de pressão de núcleo quente.
Aprofundamento: vento gradiente de perfis dados; surge com plataforma
dada. Síntese: ``por que não há furacões no Atlântico Sul — e o que o
Catarina ensinou?''. \emph{Confira a numeração na sua cópia (v1.02b).}

\subsection*{Uso do notebook \texttt{cap16\_furacoes.ipynb}}

\begin{center}
\begin{tabular}{lll}
\hline
Caso & Conteúdo & Momento sugerido \\
\hline
1 & Carnot + Clausius--Clapeyron: MPI(TSM) & Aula 1 \\
2 & Holland: $P(r)$ e $v(r)$ por categoria & Aula 2 \\
3 & \texttt{tcpyPI}: MPI de sondagem real & Aula 1, fecho \\
4 & surge: plataforma rasa × profunda & Aula 2 \\
\hline
\end{tabular}
\end{center}

O Caso~3 usa o pacote de pesquisa real (Bister--Emanuel) — vale
mostrar que a ferramenta de sala é a mesma dos artigos.

\subsection*{Exercícios complementares e soluções}

Nas notas: T1--T5 e N1--N4. Soluções em \texttt{cap16\_solucoes.ipynb}
(\emph{não distribuir}): T2 com \texttt{sympy}; N2 ajusta Holland por
otimização; N3 é a ponte quantitativa para o Cap.~21. Pares:
T1$\leftrightarrow$N1, T2$\leftrightarrow$N2, T4$\leftrightarrow$N4.

\subsection*{Conexões com o resto do curso}

Clausius--Clapeyron e calor latente (Cap.~4); fluxos ar--mar
(Cap.~3); vento gradiente (Cap.~10); alísios, $\beta$ e recurvamento
(Cap.~11); convecção da parede (Cap.~14); MPI num clima quente
(Cap.~21); previsão de trajetória por ensemble (Cap.~20).
